\documentclass[10pt,a4paper]{amsart}
\usepackage{amsmath,amssymb,amsthm,graphicx}
\usepackage{verbatim}
\newenvironment{reference}{\comment}{\endcomment}
%\newenvironment{reference}{}{}
\newenvironment{slogan}{\comment}{\endcomment}
\newenvironment{history}{\comment}{\endcomment}
% For commutative diagrams we use Xy-pic
\usepackage[all]{xy}

% We use 2cell for 2-commutative diagrams.
\xyoption{2cell}
% Theorem environments.
%
\theoremstyle{plain}
\newtheorem{theorem}[subsection]{Theorem}
\newtheorem{proposition}[subsection]{Proposition}
\newtheorem{lemma}[subsection]{Lemma}

\theoremstyle{definition}
\newtheorem{definition}[subsection]{Definition}
\newtheorem{example}[subsection]{Example}
\newtheorem{exercise}[subsection]{Exercise}
\newtheorem{situation}[subsection]{Situation}

\theoremstyle{remark}
\newtheorem{remark}[subsection]{Remark}
\newtheorem{remarks}[subsection]{Remarks}

\numberwithin{equation}{subsection}

% Macros
%
\def\lim{\mathop{\mathrm{lim}}\nolimits}
\def\colim{\mathop{\mathrm{colim}}\nolimits}
\def\Spec{\mathop{\mathrm{Spec}}}
\def\Hom{\mathop{\mathrm{Hom}}\nolimits}
\def\Ext{\mathop{\mathrm{Ext}}\nolimits}
\def\SheafHom{\mathop{\mathcal{H}\!\mathit{om}}\nolimits}
\def\SheafExt{\mathop{\mathcal{E}\!\mathit{xt}}\nolimits}
\def\Sch{\mathit{Sch}}
\def\Mor{\mathop{\mathrm{Mor}}\nolimits}
\def\Ob{\mathop{\mathrm{Ob}}\nolimits}
\def\Sh{\mathop{\mathit{Sh}}\nolimits}
\def\NL{\mathop{N\!L}\nolimits}
\def\CH{\mathop{\mathrm{CH}}\nolimits}
\def\proetale{{pro\text{-}\acute{e}tale}}
\def\etale{{\acute{e}tale}}
\def\QCoh{\mathit{QCoh}}
\def\Ker{\mathop{\mathrm{Ker}}}
\def\Im{\mathop{\mathrm{Im}}}
\def\Coker{\mathop{\mathrm{Coker}}}
\def\Coim{\mathop{\mathrm{Coim}}}

% Boxtimes
%
\DeclareMathSymbol{\boxtimes}{\mathbin}{AMSa}{"02}

%
% Macros for moduli stacks/spaces
%
\def\QCohstack{\mathcal{QC}\!\mathit{oh}}
\def\Cohstack{\mathcal{C}\!\mathit{oh}}
\def\Spacesstack{\mathcal{S}\!\mathit{paces}}
\def\Quotfunctor{\mathrm{Quot}}
\def\Hilbfunctor{\mathrm{Hilb}}
\def\Curvesstack{\mathcal{C}\!\mathit{urves}}
\def\Polarizedstack{\mathcal{P}\!\mathit{olarized}}
\def\Complexesstack{\mathcal{C}\!\mathit{omplexes}}
% \Pic is the operator that assigns to X its picard group, usage \Pic(X)
% \Picardstack_{X/B} denotes the Picard stack of X over B
% \Picardfunctor_{X/B} denotes the Picard functor of X over B
\def\Pic{\mathop{\mathrm{Pic}}\nolimits}
\def\Picardstack{\mathcal{P}\!\mathit{ic}}
\def\Picardfunctor{\mathrm{Pic}}
\def\Deformationcategory{\mathcal{D}\!\mathit{ef}}
\usepackage{cleveref}
\crefname{originalreference}{original reference}{original references}
\Crefname{originalreference}{Original reference}{Original references}
\makeatletter
\newlabel{decent-spaces-proposition-reasonable-sober}{{decent-spaces-proposition-reasonable-sober}{source}{}{sourcepointer-1}{}}
\newlabel{decent-spaces-proposition-reasonable-sober@cref}{{[originalreference][0][]decent-spaces-proposition-reasonable-sober}{[1][0][]source}{}{}{}}
\makeatother
\crefname{originallemma}{original lemma}{original lemmas}
\Crefname{originallemma}{Original lemma}{Original lemmas}
\makeatletter
\newlabel{lemma-base-change-quasi-compact}{{lemma-base-change-quasi-compact}{source}{}{sourcepointer-2}{}}
\newlabel{lemma-base-change-quasi-compact@cref}{{[originallemma][0][]lemma-base-change-quasi-compact}{[1][0][]source}{}{}{}}
\makeatother
\crefname{originallemma}{original lemma}{original lemmas}
\Crefname{originallemma}{Original lemma}{Original lemmas}
\makeatletter
\newlabel{lemma-gerbe-bijection-points}{{lemma-gerbe-bijection-points}{source}{}{sourcepointer-3}{}}
\newlabel{lemma-gerbe-bijection-points@cref}{{[originallemma][0][]lemma-gerbe-bijection-points}{[1][0][]source}{}{}{}}
\makeatother
\crefname{originallemma}{original lemma}{original lemmas}
\Crefname{originallemma}{Original lemma}{Original lemmas}
\makeatletter
\newlabel{lemma-gerbe-decent}{{lemma-gerbe-decent}{source}{}{sourcepointer-4}{}}
\newlabel{lemma-gerbe-decent@cref}{{[originallemma][0][]lemma-gerbe-decent}{[1][0][]source}{}{}{}}
\makeatother
\crefname{originallemma}{original lemma}{original lemmas}
\Crefname{originallemma}{Original lemma}{Original lemmas}
\makeatletter
\newlabel{lemma-kolmogorov}{{lemma-kolmogorov}{source}{}{sourcepointer-5}{}}
\newlabel{lemma-kolmogorov@cref}{{[originallemma][0][]lemma-kolmogorov}{[1][0][]source}{}{}{}}
\makeatother
\crefname{originallemma}{original lemma}{original lemmas}
\Crefname{originallemma}{Original lemma}{Original lemmas}
\makeatletter
\newlabel{lemma-monomorphism-cartesian-square-inertia}{{lemma-monomorphism-cartesian-square-inertia}{source}{}{sourcepointer-6}{}}
\newlabel{lemma-monomorphism-cartesian-square-inertia@cref}{{[originallemma][0][]lemma-monomorphism-cartesian-square-inertia}{[1][0][]source}{}{}{}}
\makeatother
\crefname{originallemma}{original lemma}{original lemmas}
\Crefname{originallemma}{Original lemma}{Original lemmas}
\makeatletter
\newlabel{lemma-representable-decent}{{lemma-representable-decent}{source}{}{sourcepointer-7}{}}
\newlabel{lemma-representable-decent@cref}{{[originallemma][0][]lemma-representable-decent}{[1][0][]source}{}{}{}}
\makeatother
\crefname{originalproposition}{original proposition}{original propositions}
\Crefname{originalproposition}{Original proposition}{Original propositions}
\makeatletter
\newlabel{proposition-open-stratum}{{proposition-open-stratum}{source}{}{sourcepointer-8}{}}
\newlabel{proposition-open-stratum@cref}{{[originalproposition][0][]proposition-open-stratum}{[1][0][]source}{}{}{}}
\makeatother
\crefname{originalreference}{original reference}{original references}
\Crefname{originalreference}{Original reference}{Original references}
\makeatletter
\newlabel{stacks-properties-definition-reduced-induced-stack}{{stacks-properties-definition-reduced-induced-stack}{source}{}{sourcepointer-9}{}}
\newlabel{stacks-properties-definition-reduced-induced-stack@cref}{{[originalreference][0][]stacks-properties-definition-reduced-induced-stack}{[1][0][]source}{}{}{}}
\makeatother
\title{The space of a decent stack with quasi-compact inertia is sober}
\author{The Stacks Project Authors}
\date{stacks-project:stacks-morphisms.tex:local-snapshot-20261012; GNU FDL 1.2 or later}
\begin{document}
\maketitle
\noindent\textit{Editorial note.} All delivered mathematical text is the author's own, copied from the pinned Stacks Project chapter source sources/stacks-algebraic-stacks-morphisms/pinned-original-source.tex; no statement or proof step was rewritten, reordered or removed.\par\medskip
\noindent\textit{Editorial note.} Preamble provenance: the Stacks Project keeps one single shared preamble for every chapter; this card takes it directly from the staged sources/stacks-algebraic-stacks-morphisms/preamble.tex as three verbatim spans (lines 8-12, 14-18 and 154-217), so no derived file was written and nothing was added or redefined. The wrapper class and the base packages are supplied by the assembler.\par\medskip
\noindent\textit{Editorial note.} Bibliography provenance: neither the counted statement nor its proof contains a citation (this lane rescans the whole span before writing the card), so the card carries no bibliography and no entry was reconstructed.\par\medskip
\noindent\textit{Editorial note.} Selection provenance: this card was opened under the file-based control contract (CONTROL.md), with the mechanical pre-card gate run on its statement and proof interval before the card was written.\par\medskip
\noindent\textit{Allowed standard prerequisites (editorial).}\begin{itemize}
\item The reduced induced closed substack structure on an irreducible closed subset. Invoked inside the counted proof (source lines 10584-10587) as a cross-chapter pointer. Turns the given irreducible closed subset into a stack to which the induction hypothesis applies.
\item The permanence of decency under representable morphisms and the base change description of the inertia. Invoked inside the counted proof (source lines 10588-10597). Shows that the reduced closed substack again satisfies the hypotheses.
\item The existence of a dense open substack that is a gerbe. Invoked inside the counted proof (source lines 10602-10604). Reduces further to the case of a gerbe.
\item The bijection of the space of a gerbe with the space of its base and the sobriety of decent algebraic spaces. Invoked near the end of the counted proof (source lines 10605-10616). Identifies the space of the gerbe with a decent algebraic space and applies sobriety there.
\end{itemize}\medskip
\ifdefined\hypertarget\hypertarget{sourcepointer-1}{}\fi
\noindent\textit{Source reference.} Original reference decent-spaces-proposition-reasonable-sober: sources/stacks-algebraic-stacks-morphisms/pinned-original-source.tex, source lines 10616--10616. Navigation pointer only: the counted text refers to the Stacks item decent-spaces-proposition-reasonable-sober, which belongs to another chapter of the project and is not part of this card. The pointer carries the target's own label as its number because this chapter's staged snapshot does not contain the other chapter's tag table; its content is not reproduced here.\par\medskip
\ifdefined\hypertarget\hypertarget{sourcepointer-2}{}\fi
\noindent\textit{Source reference.} Original lemma lemma-base-change-quasi-compact: sources/stacks-algebraic-stacks-morphisms/pinned-original-source.tex, source lines 1984--1984. Navigation pointer only: the counted text refers to the Stacks item lemma-base-change-quasi-compact, which is an existing labeled result of this same chapter and is not part of this card. The pointer carries the target's own label as its number and its definition line as its span; its content is not reproduced here.\par\medskip
\ifdefined\hypertarget\hypertarget{sourcepointer-3}{}\fi
\noindent\textit{Source reference.} Original lemma lemma-gerbe-bijection-points: sources/stacks-algebraic-stacks-morphisms/pinned-original-source.tex, source lines 6618--6618. Navigation pointer only: the counted text refers to the Stacks item lemma-gerbe-bijection-points, which is an existing labeled result of this same chapter and is not part of this card. The pointer carries the target's own label as its number and its definition line as its span; its content is not reproduced here.\par\medskip
\ifdefined\hypertarget\hypertarget{sourcepointer-4}{}\fi
\noindent\textit{Source reference.} Original lemma lemma-gerbe-decent: sources/stacks-algebraic-stacks-morphisms/pinned-original-source.tex, source lines 10493--10493. Navigation pointer only: the counted text refers to the Stacks item lemma-gerbe-decent, which is an existing labeled result of this same chapter and is not part of this card. The pointer carries the target's own label as its number and its definition line as its span; its content is not reproduced here.\par\medskip
\ifdefined\hypertarget\hypertarget{sourcepointer-5}{}\fi
\noindent\textit{Source reference.} Original lemma lemma-kolmogorov: sources/stacks-algebraic-stacks-morphisms/pinned-original-source.tex, source lines 10519--10519. Navigation pointer only: the counted text refers to the Stacks item lemma-kolmogorov, which is an existing labeled result of this same chapter and is not part of this card. The pointer carries the target's own label as its number and its definition line as its span; its content is not reproduced here.\par\medskip
\ifdefined\hypertarget\hypertarget{sourcepointer-6}{}\fi
\noindent\textit{Source reference.} Original lemma lemma-monomorphism-cartesian-square-inertia: sources/stacks-algebraic-stacks-morphisms/pinned-original-source.tex, source lines 1480--1480. Navigation pointer only: the counted text refers to the Stacks item lemma-monomorphism-cartesian-square-inertia, which is an existing labeled result of this same chapter and is not part of this card. The pointer carries the target's own label as its number and its definition line as its span; its content is not reproduced here.\par\medskip
\ifdefined\hypertarget\hypertarget{sourcepointer-7}{}\fi
\noindent\textit{Source reference.} Original lemma lemma-representable-decent: sources/stacks-algebraic-stacks-morphisms/pinned-original-source.tex, source lines 10425--10425. Navigation pointer only: the counted text refers to the Stacks item lemma-representable-decent, which is an existing labeled result of this same chapter and is not part of this card. The pointer carries the target's own label as its number and its definition line as its span; its content is not reproduced here.\par\medskip
\ifdefined\hypertarget\hypertarget{sourcepointer-8}{}\fi
\noindent\textit{Source reference.} Original proposition proposition-open-stratum: sources/stacks-algebraic-stacks-morphisms/pinned-original-source.tex, source lines 6746--6746. Navigation pointer only: the counted text refers to the Stacks item proposition-open-stratum, which is an existing labeled result of this same chapter and is not part of this card. The pointer carries the target's own label as its number and its definition line as its span; its content is not reproduced here.\par\medskip
\ifdefined\hypertarget\hypertarget{sourcepointer-9}{}\fi
\noindent\textit{Source reference.} Original reference stacks-properties-definition-reduced-induced-stack: sources/stacks-algebraic-stacks-morphisms/pinned-original-source.tex, source lines 6968--6968. Navigation pointer only: the counted text refers to the Stacks item stacks-properties-definition-reduced-induced-stack, which belongs to another chapter of the project and is not part of this card. The pointer carries the target's own label as its number because this chapter's staged snapshot does not contain the other chapter's tag table; its content is not reproduced here.\par\medskip
\begin{proposition}
\label{proposition-decent-sober}
Let $\mathcal{X}$ be a decent algebraic stack such that
$\mathcal{I}_\mathcal{X} \to \mathcal{X}$ is quasi-compact.
Then $|\mathcal{X}|$ is sober.
\end{proposition}

\begin{proof}
By Lemma \ref{lemma-kolmogorov} we know that $|\mathcal{X}|$ is Kolmogorov
(in fact we will reprove this).
Let $T \subset |\mathcal{X}|$ be an irreducible closed subset.
We have to show $T$ has a generic point.
Let $\mathcal{Z} \subset \mathcal{X}$ be the reduced induced
closed substack corresponding to $T$, see Properties of Stacks, Definition
\ref{stacks-properties-definition-reduced-induced-stack}.
Since $\mathcal{Z} \to \mathcal{X}$ is a closed immersion,
we see that $\mathcal{Z}$ is a decent algebraic stack, see
Lemma \ref{lemma-representable-decent}.
Also, the morphism $\mathcal{I}_\mathcal{Z} \to \mathcal{Z}$
is the base change of $\mathcal{I}_\mathcal{X} \to \mathcal{X}$
(Lemma \ref{lemma-monomorphism-cartesian-square-inertia}).
Hence $\mathcal{I}_\mathcal{Z} \to \mathcal{Z}$ is quasi-compact
(Lemma \ref{lemma-base-change-quasi-compact}).
Thus we reduce to the case discussed in the next paragraph.

\medskip\noindent
Assume $\mathcal{X}$ is decent,
$\mathcal{I}_\mathcal{X} \to \mathcal{X}$ is quasi-compact,
$\mathcal{X}$ is reduced, and $|\mathcal{X}|$ irreducible.
We have to show $|\mathcal{X}|$ has a generic point.
By Proposition \ref{proposition-open-stratum}.
there exists a dense open substack $\mathcal{U} \subset \mathcal{X}$
which is a gerbe. In other words, $|\mathcal{U}| \subset |\mathcal{X}|$
is open dense. Thus we may assume that $\mathcal{X}$ is a gerbe
in addition to all the other properties.
Say $\mathcal{X} \to X$ turns $\mathcal{X}$ into a gerbe over the
algebraic space $X$. Then $|\mathcal{X}| \cong |X|$ by
Lemma \ref{lemma-gerbe-bijection-points}.
In particular, $X$ is quasi-compact and $|X|$ is irreducible.
Also, by Lemma \ref{lemma-gerbe-decent}
we see that $X$ is a decent algebraic space.
Then $|\mathcal{X}| = |X|$ is sober by
Decent Spaces, Proposition \ref{decent-spaces-proposition-reasonable-sober}
and hence has a (unique) generic point.
\end{proof}

\end{document}
